\section{Conclusion}\label{sec:conclusion}
Under independent sub-Gaussian factors, the dual covariance separates into a finite-rank score term, an isotropic bulk shift, and a controlled remainder.  Sequential NRM and bNRM estimate the shift from one Gram matrix; CDM prevents the shift from appearing by crossing independent sample halves.  Their bulk errors depend on the same Frobenius and operator scales, but their score errors differ because CDM splits the sample.

Fixed-$n$ and growing-$n$ limits must not be conflated.  At fixed $n$, the natural targets are random score oracles.  Population eigenvalue consistency becomes available when the score covariance also concentrates.  Sequential NRM has an additional later-spike subtraction, while oracle-rank bNRM avoids it at the price of knowing the signal rank.  Finally, directions require a stronger signal than eigenvalues because they must overcome the full bulk energy $\tau_0$, not only the quadratic fluctuation $q_0$.  The numerical results reproduce all four distinctions with recorded seeds and code included in the source package.
